\section{Introduction}\label{sec:intro}
\begin{figure}[t]
\centering
\begin{tikzpicture}[
  font=\small,
  box/.style={draw, rounded corners=2pt, align=center, minimum height=9mm, inner sep=4pt, fill=white},
  stage/.style={box, minimum width=30mm},
  store/.style={box, fill=gray!10, minimum width=26mm},
  arr/.style={-{Latex[length=2mm]}, thick},
  lbl/.style={font=\scriptsize, fill=white, inner sep=1pt}]
  \node[stage] (charge) {(i) Charger\\assignment, Alg.~\ref{alg:charge}};
  \node[stage, right=9mm of charge] (busy) {(ii) Plan busy robots\\Alg.~\ref{alg:pp}};
  \node[stage, right=9mm of busy] (alloc) {(iii) Price idle robots\\$A^{\mathcal{B}}(p^r,p_\tau,t_0)$, Hungarian};
  \node[stage, right=9mm of alloc] (newp) {(iv) Plan assigned\\and parking robots};
  \node[store, below=9mm of busy, xshift=22mm] (table) {Reservation table $\mathcal{B}$\\vertices, edges, holds on $[t_0,t_0+w]$};
  \node[store, left=9mm of table] (tasks) {Pending tasks\\endpoint\\capacities $\kappa_{\mathrm{d}},\kappa_{\mathrm{p}}$};
  \node[store, right=9mm of table] (exec) {Execute $h$ ticks\\conflict check, battery, events};
  \draw[arr] (charge) -- (busy);
  \draw[arr] (busy) -- (alloc);
  \draw[arr] (alloc) -- (newp);
  \draw[arr] (busy.south) -- node[lbl, pos=0.7] {write} (busy.south |- table.north);
  \draw[arr] (table.north) -- node[lbl] {read} (alloc.south);
  \draw[arr] (newp.south) |- node[lbl, pos=0.75] {write} ([xshift=-8mm, yshift=3mm]table.north east) -- ([xshift=-8mm]table.north east);
  \draw[arr, preaction={draw=white, line width=3pt, -}] (tasks.north) |- ++(0,4mm) -| node[lbl, pos=0.15] {pool} (alloc.south west);
  \draw[arr] (table.east) -- (exec.west);
  \draw[arr, dashed] (exec.south) -- ++(0,-5mm) -| ([xshift=-6mm]charge.south) node[lbl, pos=0.25] {next epoch ($h$ ticks, idle robot, depletion)};
\end{tikzpicture}
\caption{\textbf{One replanning epoch of \sname.} Static robots are written as holds and low-charge robots are matched to free charger slots; busy robots, including those just sent to a charger, are then planned into a fresh reservation table; every idle robot is then priced against the candidate pool by running the same windowed space-time A* against the table (a reservation-aware pickup cost), the Hungarian method assigns tasks, and the assigned and parking robots are planned last. A hold cascade (Alg.~\ref{alg:pp}) keeps the table conflict-free whenever a robot cannot find a path.}
\label{fig:overview}
\end{figure}


An operator console for a fleet of autonomous mobile robots (AMRs) shows the
same few things on every screen: which robot is moving where, which task it
carries, how much battery it has and which zone or dock it is heading for.
Behind those screens sits a fleet management system (FMS) that decides all of
it, continuously: new tasks arrive, idle robots must be matched to them,
every robot needs a path that does not collide with any other, and robots
must be sent to chargers before they strand. Each of these problems is well
studied on its own---multi-robot task allocation
\citep{gerkey2004formal,korsah2013comprehensive}, multi-agent path finding
(MAPF) \citep{stern2019mapf,sharon2015cbs} and its lifelong
pickup-and-delivery variant \citep{ma2017lifelong,li2021rhcr}---but a
deployable FMS has to run them together, at every tick, with hard safety
constraints.

The interface between allocation and path finding is where most fleet
managers cut corners. Assignment costs are almost always free-space distances
\citep{ma2017lifelong,liu2019task}, which ignore the traffic that the planner
already knows about; the methods that couple assignment and MAPF exactly
(CBM, \citealp{ma2016optimal}; CBS-TA, \citealp{honig2018cbsta}) solve
one-shot instances and do not transfer to a rolling-horizon loop with
batteries, dwell times and dock queues, and the lifelong scheme that does use
planned arrival times, token passing with task swaps \citep{ma2017lifelong},
applies them greedily, one robot at a time. Whether the cheap decoupling actually
costs anything in a lifelong warehouse setting is, to our knowledge, an open
empirical question; and the planners that make lifelong operation scale
\citep{li2021rhcr} are evaluated without batteries and without dwell times
or capacity limits at the delivery endpoints.

To answer that question with a system rather than a thought experiment, we
introduce \emph{\sname} (Path-Aware Coupled Tasking), an end-to-end,
framework-free lifelong fleet manager for grid warehouses
(\cref{fig:overview}). \sname\ replans the fleet every $h$ ticks over a window
of $w$ ticks, as in RHCR \citep{li2021rhcr}, with windowed space-time A*
\citep{silver2005cooperative} under a reservation table and prioritized
planning \citep{erdmann1987multiple}. Its coupling is deliberately the
cheapest one that uses real plans: the cost of assigning robot $r$ to task
$\tau$ is the arrival time that the same planner reports for $r$'s trip to
the pickup against the reservations of the robots that are already busy, and
the Hungarian method \citep{kuhn1955hungarian} turns those costs into an
optimal assignment. Around the coupling sit the pieces that make a lifelong
loop robust: a hold cascade that converts planner failures into consistent
rests, token-passing endpoint capacities \citep{ma2017lifelong} that keep dock
queues short, a two-threshold charging policy with capacity-constrained slot
assignment, and an optional windowed Conflict-Based Search backend
\citep{sharon2015cbs}.

Our contributions are:
\begin{itemize}
  \item \emph{A tested lifelong FMS stack}: grid map model with benchmark
  layouts (one generated from an operator console's own zone and dock
  vocabulary), Hungarian and greedy allocation, windowed space-time A*,
  prioritized planning with a hold cascade, windowed CBS, charging, and a
  seeded discrete-time simulator; conflict-freeness of every executed step
  and optimality of the assignment step are proved, and a test suite checks the
  Hungarian and CBS implementations against exhaustive search and the planners
  against property tests.
  \item \emph{Path-aware coupled allocation}: assignment costs from
  MAPF-in-the-loop arrival times, with a congestion-proxy middle ground, a
  battery term and candidate pruning, designed so that the coupling adds only
  a bounded number of windowed A* calls per epoch.
  \item \emph{A controlled study of what moves a fleet}: 420 seeded runs on
  three layouts sweep allocation cost model, planner backend, fleet size,
  load regime, window, replanning period, dock capacity and battery
  parameters; a second round of 6{,}180 runs adds a 60-seed safety stress
  test checked by an independent trajectory auditor, an ablation of the hold
  cascade and of the priorities, token-passing baselines and four regimes
  chosen to favour the coupling and committed before their results. Every number traces to a
  committed result file.
\end{itemize}

The study's main result is negative and, we argue, useful: in every one of
twelve layout/fleet/load cells the four allocation cost models---greedy,
Hungarian, congestion proxy and MAPF-in-the-loop---differ by at most 2.9\% in
mean throughput, and 35 of 36 seed-paired 95\% intervals against Hungarian
contain zero (\cref{tab:main,tab:paired}). With five seeds this rules out
differences above about 7\% in the noisiest cell and above about 4\% under the Poisson
load, not small ones. What does move the fleet is elsewhere: raising the
per-dock in-flight capacity from 2 to 3 raises throughput by 26--29\%,
shortening the planning window from 20 to 10 ticks raises it by 7\% while
cutting planner time by $3.3\times$, and windowed CBS reduces makespan by
8\% on small fleets (\cref{tab:ablation,tab:planner}); all three effects lie
outside their seed-paired 95\% intervals. The negative result holds in the
four regimes we chose to favour the coupling (\cref{subsec:coupling}).
Token passing \citep{ma2017lifelong} is 16--66\% slower than our loop at
its default dock capacity, but its endpoint rule allows one robot in flight
per dock; at that capacity our loop and token passing with task swaps are
within $-2.5$ to $+1.6\%$ of each other (\cref{subsec:baselines}), so the
gap measures dock capacity, not our allocation or planner.
Throughout, the executed plans of 60.4 million audited robot-steps contain no
vertex or swap conflict, as the construction guarantees. Progress and charge
are not guaranteed: with normal batteries no robot depletes and every run
delivers all of its tasks, but with doubled drain the charging policy
strands robots once its slots saturate (\cref{subsec:main}), and we prove
only a conditional reserve bound (\cref{prop:charge}). We release the
implementation inside the operator console it was built for, together with
all results.
